\section{Criterio de selección}
\label{sec:candidate-selection}

\HAFOCardCandidates

La selección cruzó el catálogo de sistemas, los manifiestos de validación, el
manuscrito metodológico y las fuentes bibliográficas primarias. Se priorizaron
sistemas que cumplen simultáneamente:

\begin{enumerate}
 \item derivada ordinaria o Caputo conmensurable \(q\in(0,1]\);
 \item ecuaciones, parámetros, orden e iniciales publicados;
 \item no linealidades representables por D0, T1, A1 o M1;
 \item un reclamo publicado de atractor oculto, cuenca o coexistencia;
 \item tamaño compatible con ABM de memoria completa y una reproducción
 independiente;
 \item valor científico como prueba de una parte aún no cerrada del método.
\end{enumerate}

Se excluyen del carril actual sistemas incommensurables, Caputo--Fabrizio,
Atangana--Baleanu, retardados, discretos o complejos. No son imposibles, pero
su símbolo de transferencia, su condición de estabilidad y su contrato de
memoria no son intercambiables con Caputo conmensurable.

\section{Prioridad 1: sistema fraccionario sin equilibrio de Muñoz--Pacheco}
\label{sec:cand-munoz}

La fuente primaria \cite{munoz2018} propone
\begin{equation}
\begin{aligned}
 {}^{\mathrm C}D^q x&=yz+x(y-a),\\
 {}^{\mathrm C}D^q y&=1-|x|,\\
 {}^{\mathrm C}D^q z&=-xy-z.
\end{aligned}
\label{eq:cand-munoz}
\end{equation}
Para \(a=0.35,q=0.97\), publica un atractor caótico desde
\((1,1,1)\). Para \(a=0.35,q=0.996\), publica coexistencia de un candidato
caótico desde \((1,1,1)\) y uno periódico desde \((0,75,-50)\), junto con un
corte de cuenca. Esto lo convierte en el candidato fraccionario no Chua más
completo encontrado.

\subsection{Realización M1 exacta}

Con \(X=(x,y,z)^{\mathsf T}\),
\[
 \Phi(X)=(yz,xy,|x|)^{\mathsf T},
\]
se tiene
\begin{equation}
 {}^{\mathrm C}D^qX=A X+d+B\Phi(X),
 \label{eq:cand-munoz-mimo}
\end{equation}
donde
\begin{equation}
 A=\begin{pmatrix}
 -a&0&0\\0&0&0\\0&0&-1
 \end{pmatrix},\quad
 d=\begin{pmatrix}0\\1\\0\end{pmatrix},\quad
 B=\begin{pmatrix}
 1&1&0\\
 0&0&-1\\
 0&-1&0
 \end{pmatrix}.
 \label{eq:cand-munoz-matrices}
\end{equation}
La multiplicación da
\[
 AX+d+B\Phi=
 \begin{pmatrix}
 -ax+yz+xy\\1-|x|\\-z-xy
 \end{pmatrix},
\]
idéntica a \cref{eq:cand-munoz}. Wolfram devolvió el residual
\(\{0,0,0\}\).

Con \(C=I\) y \(z_f=s^q\), la matriz de transferencia es especialmente simple:
\begin{equation}
 W(z_f)=(z_fI-A)^{-1}B
 =\begin{pmatrix}
 \dfrac1{z_f+a}&\dfrac1{z_f+a}&0\\[0.8em]
 0&0&-\dfrac1{z_f}\\[0.8em]
 0&-\dfrac1{z_f+1}&0
 \end{pmatrix}.
 \label{eq:cand-munoz-transfer}
\end{equation}
El término constante \(d\) obliga a resolver el balance DC; los productos
cuadráticos generan segundo armónico y \(|x|\) exige cuadratura PWA exacta o
una homotopía suavizada. El cierre mínimo es M1 con \(H\ge2\).

\subsection{Demostración completa de ausencia de equilibrios}

En equilibrio, la segunda ecuación exige \(|x|=1\), de modo que
\(x=\pm1\neq0\). La tercera da
\[
 z=-xy.
\]
Al sustituir en la primera,
\[
 0=y(-xy)+x(y-a)
 =x(-y^2+y-a).
\]
Como \(x\neq0\), debe cumplirse
\begin{equation}
 y^2-y+a=0.
 \label{eq:cand-munoz-equilibrium}
\end{equation}
Su discriminante es
\[
 \Delta_y=1-4a.
\]
Para \(a=0.35\), \(\Delta_y=-0.4<0\); por tanto, no existe equilibrio real.
Wolfram confirmó \texttt{Reduce[...,Reals]=False}. Bajo la definición usual,
cualquier atractor del sistema es oculto por ausencia de equilibrios. Aun así,
deben reproducirse recurrencia y caos: la ausencia de equilibrios no transforma
una trayectoria transitoria o divergente en atractor.

Para \(x\neq0\), el Jacobiano regional es
\begin{equation}
 J(X)=
 \begin{pmatrix}
 y-a&z+x&y\\
 -\operatorname{sign}(x)&0&0\\
 -y&-x&-1
 \end{pmatrix}.
 \label{eq:cand-munoz-jacobian}
\end{equation}
No se aplica Matignon en equilibrios porque éstos no existen. La validación
debe concentrarse en el integrador nonsuave, la coexistencia y la robustez de
los diagnósticos.

\begin{proposalbox}
\textbf{Plan L0--L6.} Congelar \(a=0.35\), \(q=0.97\) y \(q=0.996\), las dos
iniciales y el convenio de \(|x|\); derivar coeficientes M1; resolver \(H=2,3,4\)
y controlar \(\eta_H\); continuar una semilla; integrar Caputo con memoria
completa; reproducir la cuenca \(x=0\); comparar ABM con otro integrador. Como
no hay equilibrios, L5 se sustituye por exploración de cuenca, coexistencia y
frontera de divergencia.
\end{proposalbox}

\section{Prioridad 2: jerk cuadrático entero con equilibrio estable}
\label{sec:cand-jerk}

Liu, Sang, Wang y Ahmad \cite{liu2021} estudian
\begin{equation}
 \dot x=y,\qquad
 \dot y=z,\qquad
 \dot z=-x-y-1.1z-2.59z^2+xy.
 \label{eq:cand-jerk}
\end{equation}
La fuente reporta un atractor caótico oculto desde
\((0,-1.5,0)\), además de cortes de cuenca y una realización de circuito.

\subsection{Parte lineal, equilibrio y estabilidad}

Con \(X=(x,y,z)^{\mathsf T}\),
\begin{equation}
 A=\begin{pmatrix}
 0&1&0\\0&0&1\\-1&-1&-1.1
 \end{pmatrix}.
 \label{eq:cand-jerk-A}
\end{equation}
Las dos primeras ecuaciones de equilibrio dan \(y=z=0\); la tercera da
\(x=0\). El origen es el único equilibrio. Su polinomio característico es
\begin{equation}
 p(s)=s^3+1.1s^2+s+1.
 \label{eq:cand-jerk-charpoly}
\end{equation}
El criterio de Routh--Hurwitz para un cúbico exige coeficientes positivos y
\(1.1\cdot1>1\), ambas condiciones satisfechas. El equilibrio único es
asintóticamente estable. Un atractor coexistente no puede localizarse
siguiendo la variedad inestable del origen y, si su cuenca está separada de una
vecindad del mismo, es oculto.

\subsection{Realización M1 y obstrucción escalar}

Definimos
\[
 \Phi(X)=(xy,z^2)^{\mathsf T},\qquad
 B=\begin{pmatrix}0&0\\0&0\\1&-2.59\end{pmatrix}.
\]
Entonces
\[
 AX+B\Phi(X)=
 (y,z,-x-y-1.1z+xy-2.59z^2)^{\mathsf T},
\]
y Wolfram devuelve residual nulo.

Aunque los dos canales actúan en la dirección \(e_3\), no pueden reducirse a
una no linealidad cuadrática de un único observable. Para
\[
 Q(x,y,z)=xy-2.59z^2
\]
se tiene
\begin{equation}
 \nabla^2Q=
 \begin{pmatrix}
 0&1&0\\1&0&0\\0&0&-5.18
 \end{pmatrix},
 \qquad \operatorname{rank}\nabla^2Q=3.
 \label{eq:cand-jerk-hessian}
\end{equation}
En cambio, \(Q=\kappa(r^{\mathsf T}X)^2\) implicaría
\(\nabla^2Q=2\kappa rr^{\mathsf T}\), de rango a lo sumo uno. Esto falsifica
D0 y justifica M1.

Como ambas entradas llegan al último estado, la transferencia es
\begin{equation}
 W(s)=\frac1{p(s)}
 \begin{pmatrix}
 1&-2.59\\
 s&-2.59s\\
 s^2&-2.59s^2
 \end{pmatrix}.
 \label{eq:cand-jerk-transfer}
\end{equation}
Los términos cuadráticos requieren DC y segundo armónico. Este sistema es,
por tanto, el primer banco entero recomendado para demostrar que M1 añade
capacidad real sobre los casos escalares ya cerrados.

\section{Prioridad 3: sistema Liu extendido fraccionario}
\label{sec:cand-extended-liu}

La fuente \cite{wang2020} usa
\begin{equation}
\begin{aligned}
 {}^{\mathrm C}D^q x&=a(y-x),\\
 {}^{\mathrm C}D^q y&=bx-kxz+d\,\operatorname{sign}(w)+m,\\
 {}^{\mathrm C}D^q z&=-cz+hx^2,\\
 {}^{\mathrm C}D^q w&=-gy.
\end{aligned}
\label{eq:cand-liu}
\end{equation}
Una realización exacta usa
\[
 \Phi=(xz,x^2,\operatorname{sign}w)^{\mathsf T},
\]
\[
 A=\begin{pmatrix}
 -a&a&0&0\\ b&0&0&0\\0&0&-c&0\\0&-g&0&0
 \end{pmatrix},\quad
 d_0=(0,m,0,0)^{\mathsf T},
\]
\[
 B=\begin{pmatrix}
 0&0&0\\-k&0&d\\0&h&0\\0&0&0
 \end{pmatrix}.
\]
Wolfram confirmó el residual nulo.

Para \(a=10,b=40,c=2.5,k=1,h=4,g=5,d=m=20\), las ecuaciones de equilibrio
dan sucesivamente \(y=x\), \(y=0\), \(x=0\), \(z=0\), y
\[
 20\,\operatorname{sign}(w)+20=0.
\]
Con \(\operatorname{sign}(0)=0\), la solución es la semirrecta
\begin{equation}
 x=y=z=0,\qquad w<0,
 \label{eq:cand-liu-equilibria}
\end{equation}
no un sistema sin equilibrios. Antes de promover el caso deben congelarse
\(q\), la rama dinámica, las iniciales, la convención en \(w=0\) y el protocolo
para una variedad continua de equilibrios.

\section{Cola interna y controles}
\label{sec:candidate-queue}

\begin{longtable}{@{}L{0.23\linewidth}L{0.20\linewidth}L{0.24\linewidth}L{0.25\linewidth}@{}}
\caption{Portafolio compatible y requisito de promoción.}
\label{tab:candidate-portfolio}\\
\toprule
Familia & Ruta estructural & Activo disponible & Falta decisiva\\
\midrule
\endfirsthead
\toprule
Familia & Ruta estructural & Activo disponible & Falta decisiva\\
\midrule
\endhead
Leonov--Kuznetsov (162) & D0+A1 & Ecuaciones y semilla parciales &
Congelar característica PWA y criterio por etapa.\\
Genesio--Tesi & D0 sesgado, \(H\ge2\) & Realización
\(\psi(x)=x^2\) & Parámetros, iniciales y manifiesto L0 congelados.\\
Jerk con seno hiperbólico & D0, DF con \(I_1\) & Derivación simbólica &
Resolver ambigüedad de fuente \(\sinh\) frente a función inversa.\\
Oscilador memristivo & T1 sobre variedad & Cierre cúbico sesgado &
Constante de la variedad, parámetros e iniciales.\\
Jerk PWA & D0 PWA & Estructura tipo Chua & Congelar fuente; usar como control
de generalización, no como familia distante.\\
Fischer jerk exponencial & D0 con DC & Parámetros y diagnóstico &
Control bibliográfico/Lyapunov; clasificación de cuenca no reproducida
localmente.\\
Lorenz generalizado & M1, \(H\ge2\) & Álgebra completa de este reporte &
Raíz armónica, continuación, integración y cuenca.\\
Rabinovich--Fabrikant & M1, \(H\ge3\) & Álgebra completa de este reporte &
Raíz armónica, continuación, integración y cuenca.\\
Sistema financiero Fischer & M1+PWA & Banco Lyapunov & Resolver discrepancias
y separar del modelo financiero del manuscrito.\\
Four-wing Fischer & M1 cuadrático & Banco Lyapunov & Semilla y búsqueda de
ocultamiento inexistentes.\\
Debbouche fraccionario & M1+signo & DOI y parámetros parciales
\cite{debbouche2021} & Congelar ecuaciones, órdenes e iniciales de una rama.\\
\bottomrule
\end{longtable}

\section{Orden recomendado de ejecución}
\label{sec:candidate-order}

\begin{enumerate}
 \item ejecutar primero el jerk cuadrático entero: elimina el costo de memoria
 fraccionaria y prueba de forma limpia el salto D0 \(\to\) M1;
 \item ejecutar después el sistema de Muñoz--Pacheco: prueba M1, nonsuavidad,
 Caputo y coexistencia, con ecuaciones e iniciales publicadas;
 \item resolver el balance multiharmónico de Lorenz generalizado y RF usando
 las realizaciones ya auditadas, sin prometer ocultedad hasta L5;
 \item cerrar Genesio--Tesi como control escalar sesgado;
 \item dejar Liu extendido y Debbouche para una segunda ola, porque las
 variedades de equilibrio y el signo introducen contratos adicionales.
\end{enumerate}

Esta secuencia maximiza información científica por costo: primero falsifica o
confirma la extensión M1 en ODE, luego añade memoria Caputo y nonsuavidad, y
sólo después aborda los casos bibliográficos cuya semilla no está publicada.
